\section{Experiments}
\label{sec:exp}

\paragraph{Benchmarks.}
We train on one LIBERO suite at a time~\cite{liu2023libero} and evaluate on its nominal tasks and on the LIBERO-PRO
perturbations of the same suite~\cite{zhou2025liberopro}: \emph{swap} (objects exchange places), \emph{position} (the
target is moved; LIBERO-Object only), \emph{object} (objects recolored or re-textured) and \emph{language} (instruction
rephrased). Each cell is 10 tasks $\times$ 20 initial states; horizons are 280 (Object), 220 (Spatial), 300 (Goal) and 520
(Long) steps. Seeds are 7, 8, 9 unless stated.

\paragraph{Baselines, all in the same harness, teacher and budget.}
\emph{Distill}: the on-policy distillation of \cref{sec:method} without counterfactuals.
\emph{\ect}~\cite{chen2026ect}: each iteration, 16 successful episodes are replayed in a transformed scene (robot and start
pose unchanged; Object: mirror in $y$, mirror in $x$ and in $xy$ with a shift; Spatial and Goal: mirror in $y$; Long: a
shift; furniture is reflected so that it faces the work area); a tracking controller follows the transformed end-effector
path, and successful replays are paired with the original queries, on at most 15\% of the states as for \ours.
Pairs on 50\% of the states do not help \ect (Object swap 25.0, Spatial 48.0).
\emph{\pt}~\cite{tsai2026point}: the 3D displacement from the gripper to the current sub-goal (the target until it is
lifted, then its goal region), taken from the simulator at training \emph{and} test time, is mapped by an MLP with a
zero-initialized last layer and added to the flow-time conditioning of every action-expert layer; the teacher does not see
it. Published numbers of these methods come from different training regimes (four suites jointly, 30k steps, or 69 tasks)
and are not comparable to a per-suite adaptation; we therefore compare re-implementations.

\subsection{Main results}
\label{sec:exp:main}

\begin{table*}[t]
\centering
\small
\setlength{\tabcolsep}{5pt}
\begin{tabular}{l c ccccc c cccc}
\toprule
& & \multicolumn{5}{c}{LIBERO-Object} & & \multicolumn{4}{c}{LIBERO-Spatial} \\
\cmidrule{3-7}\cmidrule{9-12}
Method & seeds & nominal & swap & position & object & language & & nominal & swap & object & language \\
\midrule
Distill (no counterfactual) & 3 / 2 & 99.3 & 18.7\,{\scriptsize$\pm$2.8} & 12.2\,{\scriptsize$\pm$1.3} & 92.5 & 99.2 & & 97.5 & 40.0\,{\scriptsize$\pm$2.8} & 98.5 & 96.5 \\
\pt~\cite{tsai2026point} (oracle point) & 1 & 98.5 & 18.5 & 12.5 & 93.0 & 100.0 & & 99.0 & 40.0 & 98.0 & 97.0 \\
\ect~\cite{chen2026ect} & 2\pending{+1} & 99.0 & 23.8\,{\scriptsize$\pm$0.4} & 16.5\,{\scriptsize$\pm$0.0} & 93.3 & 98.5 & & 98.8 & 50.8\,{\scriptsize$\pm$1.8} & 98.0 & 96.8 \\
\ours (ours) & 3 & 97.3 & \textbf{54.5}\,{\scriptsize$\pm$3.8} & \textbf{38.2}\,{\scriptsize$\pm$5.6} & 81.8\pending{$^\dagger$} & 98.0\pending{$^\dagger$} & & 98.8 & \textbf{65.0}\,{\scriptsize$\pm$1.8} & 98.8\pending{$^\dagger$} & 96.8\pending{$^\dagger$} \\
\bottomrule
\end{tabular}
\caption{\textbf{LIBERO-PRO success (\%) after per-suite adaptation of \pizf}, mean $\pm$ std over seeds (200 episodes per
cell and seed). All methods share the teacher, rollouts, budget and optimizer. The untrained \pizf scores 99 / 18.5 on
Object nominal / swap. Distill has 3 seeds on Object and 2 on Spatial. \pending{$^\dagger$: 2 seeds, third running.}}
\label{tab:main}
\end{table*}

\Cref{tab:main} shows the outcome on the two suites where the swap failure is largest.
\ours raises Object swap from 18.7 to 54.5 and Object position from 12.2 to 38.2, and Spatial swap from 40.0 to 65.0,
keeping nominal success at 97--99.
Under the same teacher and budget, \ect gains 5.1 and 10.8 points on the swap cells, \ours 35.8 and 25.0;
the difference to \ect is 30.8 points on Object swap, 21.7 on Object position and 14.3 on Spatial swap.
\pt does not move any cell.
The cost is on the \emph{object} cell of LIBERO-Object, where recolored or re-textured objects drop from 92.5 to 81.8;
we analyze this trade-off in \cref{sec:exp:limits}.

\begin{table}[t]
\centering
\small
\setlength{\tabcolsep}{3.5pt}
\begin{tabular}{lcccc}
\toprule
Counterfactual (LIBERO-Object) & seeds & swap & position & nominal \\
\midrule
none (Distill) & 3 & 18.7 & 12.2 & 99.3 \\
\multicolumn{5}{l}{\emph{whole scene transformed}} \\
\quad mirror (equivariant label) & 2 & 20.8 & 12.8 & 98.0 \\
\quad rotation about the robot & 3 & 29.2 & 16.2 & 99.2 \\
\multicolumn{5}{l}{\emph{hand--target geometry preserved}} \\
\quad hand and target co-shifted & 3 & 24.3 & 20.7 & 98.7 \\
\multicolumn{5}{l}{\emph{decision changed, privileged label}} \\
\quad retarget (before grasp) & 3 & 53.0 & 20.5 & 99.0 \\
\quad relocate (while carrying) & 2 & 22.5 & 32.8 & 98.8 \\
\quad retarget + relocate & 3 & 56.3 & 23.7 & 96.8 \\
\quad retarget/co-shift + relocate & 1 & 53.0 & 39.0 & 96.0 \\
\bottomrule
\end{tabular}
\caption{\textbf{Which counterfactual helps.} Only counterfactuals that change what the policy must do at a decision state,
labeled by the privileged teacher, raise the cells: retargeting fixes \emph{which} object (swap), relocating fixes
\emph{where} it goes (position). Without validity checks or recoloring.}
\label{tab:kinds}
\end{table}

\paragraph{Which counterfactual matters (P2--P4).}
\Cref{tab:kinds} varies the counterfactual and keeps everything else fixed.
Transforming the whole scene, by mirroring with the equivariantly mapped action or by rotating the world about the robot,
changes little, consistent with the per-suite ablations of \ect itself~\cite{chen2026ect}.
Co-shifting hand and target, which preserves the geometry the memorized action relies on, does not help either.
What helps is a counterfactual that changes the decision: moving the target before the grasp fixes the choice of object
(swap 53.0) but not the placement, and moving the container while carrying fixes the placement (position 32.8) but not the
choice; combining them fixes both.
The label matters as much as the world: on Spatial, labeling the relocated world with \pizf's own chunk instead of the
privileged placer gives swap 40.0 (= Distill), against 64.7 with the placer (3 seeds), because \pizf, as
\cref{sec:probing} shows, does not act on the relocation it sees.

\subsection{What changes inside the policy}
\label{sec:exp:mech}

\begin{table}[t]
\centering
\small
\setlength{\tabcolsep}{2.5pt}
\begin{tabular}{llcccc}
\toprule
& Probe (LIBERO-Object) & \pizf & \pt & \ect & \ours \\
\midrule
\multirow{5}{*}{\rotatebox{90}{use}} & early, target moved & .13--.22 & .12--.21 & .14--.23 & .20--.30 \\
& mid, moved 20\,cm & .31 & .32 & .37 & \textbf{.67} \\
& mid, swap & .03 & .03 & .07 & \textbf{.53} \\
& mid, target renamed & .02 & .02 & .01 & \textbf{.54} \\
& carrying, cont.\ $\geq$12\,cm & .06 & .06 & .11 & \textbf{.44--.69} \\
\midrule
\multirow{3}{*}{\rotatebox{90}{enc.}} & target moved $\geq$12\,cm & .58 & -- & .62 & .65 \\
& swap & .55 & -- & .56 & .70 \\
& container moved $\geq$12\,cm & .80 & -- & .86 & .85 \\
\bottomrule
\end{tabular}
\caption{\textbf{Use and encoding after training} (seed 7; \ours: the configuration of the last row of \cref{tab:kinds}).
Use: median follow of the action (\cref{sec:probing}); ranges span displacement sizes. Encoding: in-state probe $R^2$ of
the action tokens, late expert layers. Every model encodes the change; only \ours acts on it.}
\label{tab:mech}
\end{table}

\Cref{tab:mech} applies the probes of \cref{sec:probing} to the trained policies.
The encoding is nearly the same in the three models.
What differs is use: \ect leaves it at the level of \pizf, while \ours makes the action follow a swap (0.03 $\to$ 0.53), a
renamed target (0.02 $\to$ 0.54) and a moved container (0.06 $\to$ 0.44--0.69).
Counterfactual supervision thus does not teach the policy to see; it teaches it to act on what it sees.
One blind spot is shared by all three models: early in the approach, the action follows a moved target by only 0.13--0.30.

\paragraph{Explicit geometry is ignored (P1).}
\pt is the most direct test of whether information is missing: the policy receives the exact displacement to the
sub-goal, including after a swap.
Its use is identical to that of \pizf (\cref{tab:mech}).
The injected branch converges to a near-constant offset: its output has a norm of about 1.6 but changes by only 0.05--0.08
when the displacement changes by 10\,cm.
On fixed layouts the goal point always coincides with the memorized trajectory, so nothing forces the policy to read it.
Adding the counterfactual data of \ours to \pt (with a 10$\times$ smaller learning rate for the branch, which otherwise
destabilizes training) gives Object swap 43.0 and position 41.5, no better than \ours alone, and the branch is still
ignored ($\lVert f(\mathbf{d}) - f(\mathbf{0}) \rVert$ 0.03--0.12 for 10--30\,cm against $\lVert f(\mathbf{0})\rVert$ 1.08):
under pressure to use the target's position, the policy solves the task through its VLM rather than through the new
channel.

\subsection{Where the change lives}
\label{sec:exp:local}

\begin{table}[t]
\centering
\small
\setlength{\tabcolsep}{3pt}
\begin{tabular}{lrcccc}
\toprule
Trainable part & params & swap & pos. & nominal & object \\
\midrule
none (Distill) & -- & 18.7 & 12.2 & 99.3 & 92.5 \\
expert readout & 34K & 23.5 & 20.5 & 100 & 92.5 \\
expert, last 6 layers & 4.7M & 34.5 & 30.5 & 82.5 & 80.5 \\
action expert (2 seeds) & 13.9M & 32.0 & 29.8 & 94.3 & 87.5 \\
VLM keys and values & 2.7M & 20.0 & 9.5 & 72.0 & 63.5 \\
VLM layers 9--17 & 19.6M & 41.0 & 28.5 & 92.0 & 76.5 \\
VLM (3 seeds) & 39.2M & \textbf{49.3} & 29.8 & 96.0 & 68.7 \\
all & & 53.0 & 39.0 & 96.0 & 80.0 \\
\bottomrule
\end{tabular}
\caption{\textbf{Localization.} LIBERO-Object, the counterfactual data of the last row of \cref{tab:kinds}, only the
listed LoRA adapters trained. Adapting the VLM with the action expert frozen gives 89\% of the swap gain.}
\label{tab:local}
\end{table}

With the counterfactual data fixed, we restrict which LoRA adapters may change (\cref{tab:local}).
The readout of the action expert, the natural place for a decoding bottleneck and the layer whose re-training suffices in
image classification~\cite{kirichenko2023dfr}, gains almost nothing.
Forcing the change into the last layers of the action expert gains part of the swap cell but breaks nominal behavior
(82.5).
The full action expert recovers about 40\% of the swap gain and two thirds of the position gain.
Adapting only the VLM, with the action expert frozen, recovers 89\% of the swap gain at nominal 96, replicated over three
seeds (50.0, 44.0, 54.0); most of it comes from layers 9--17 (41.0).
Yet the change is not a re-presentation of what the action expert reads: adapting only the keys and values of the VLM, the
exact interface the expert attends to, gives nothing and damages every cell.
The fix requires the VLM to \emph{process} the scene differently in its mid-to-late layers: deciding which object to go
to is computed in the VLM, and \pizf learned to compute it from the instruction alone.
Placement (position) needs both halves.
The appearance trade-off also originates in the VLM: adapting only the VLM yields the lowest \emph{object} cell (68.7),
the action expert the highest (87.5).

\subsection{Costs, other suites, and an evaluation detail}
\label{sec:exp:limits}

\paragraph{Appearance.}
A policy that ignores the scene is robust to recoloring for free; a policy that looks at it can start to identify objects
by color.
In the VLM-only configuration, the validity checks of \cref{sec:method} raise the \emph{object} cell from 69 to 83.5;
recoloring objects during the rollouts adds only 1.5 points to it and costs 9 points of swap (48.5 $\to$ 39.5), so its
benefit is doubtful. \pending{The main model uses recoloring at probability 0.5; \ect with the same recoloring, and the main
model without it, are running.} With the oracle point of \pt added, the \emph{object} cell returns to 92.5, since the point identifies the
target regardless of its color; a non-privileged grounding source would be needed to claim this.

\paragraph{LIBERO-Goal and LIBERO-Long.}
\begin{table}[t]
\centering
\small
\setlength{\tabcolsep}{4pt}
\begin{tabular}{lcccc}
\toprule
& \multicolumn{2}{c}{Goal} & \multicolumn{2}{c}{Long} \\
\cmidrule(lr){2-3}\cmidrule(lr){4-5}
Method (seed 7) & nominal & swap & nominal & swap \\
\midrule
Distill & 97.5 & 29.0 & 95.0 & 8.5 \\
\ect & 96.5 & 24.5 & 84.0 & 14.0 \\
\ours & 89.5--94.0 & 28.5--31.0 & 84.0--87.0 & 11.0--17.0 \\
\bottomrule
\end{tabular}
\caption{\textbf{Multi-step suites.} Ranges over the \ours configurations we tried. \pending{Re-evaluation with fixed
furniture placement running.}}
\label{tab:goal_long}
\end{table}
On the multi-step suites (\cref{tab:goal_long}) neither method improves swap much, and both lose nominal success on Long.
Per-task analysis traced the Goal drop to the two drawer tasks, and, within them, to two causes.
First, \pizf opens the top drawer by pushing its handle without closing the gripper, so a gripper-based phase rule took the
drawer phase for the approach to the bowl; segmenting phases by when the rest of the scene stops changing removes these
pairs, but does not remove the drop.
Second, LIBERO \emph{re-draws the furniture placement at every reset}: the cabinet moves by up to 2\,cm in $x$ and 1.4\,cm in
$y$, and the benchmark's initial states, which store joint positions only, do not fix it.
Pinning the cabinet at controlled offsets shows that \ours, unlike \pizf and Distill, is sensitive to it on the
top-drawer task (\cref{tab:furniture}): having learned to act on where movable objects are, it lost part of \pizf's ability
to follow the furniture.
\pending{Counterfactuals that nudge the furniture and keep the teacher's own response there, which \pizf handles reliably
within this range, are running.}

\begin{table}[t]
\centering
\small
\setlength{\tabcolsep}{4pt}
\begin{tabular}{lcccc}
\toprule
Cabinet offset & \pizf & Distill & \ect & \ours \\
\midrule
none & 90 & 95 & 90 & 80 \\
$-1$\,cm in $x$ & 95 & 95 & 80 & 65 \\
$-2$\,cm in $x$ & 85 & 90 & 85 & 60 \\
$(-2, -1.4)$\,cm & 95 & 90 & 85 & 55 \\
\bottomrule
\end{tabular}
\caption{\textbf{Furniture sensitivity}, Goal task ``open the top drawer and put the bowl inside'' (20 trials, furniture
pinned at the placement of a fresh environment plus the offset; offsets within the range LIBERO itself draws). \ours: the
configuration with the best Goal nominal success in \cref{tab:goal_long} (94.0).}
\label{tab:furniture}
\end{table}

\paragraph{Evaluation protocol.}
The re-drawn furniture has a second consequence.
The standard LIBERO loop runs the trials of a task sequentially in one seeded environment, so trial $k$ sees the furniture
of the $k$-th draw, while evaluation scripts that spread trials over parallel workers give the same trial a different
placement depending on scheduling.
With one adapter, the task ``open the middle drawer'' of LIBERO-Goal scores 65--70\% when its 20 trials run sequentially
and 95\% when they run in freshly built environments, failing on exactly the same trials in repeated sequential runs.
All numbers in this paper use the same scheduler for every method; for strict comparability we seed the environment per
(task, trial) before each reset, which fixes the furniture of every trial independently of scheduling, and recommend
reporting it.
LIBERO-Object has no furniture and is unaffected.
